
\section{Related Work}

\subsection{Low-Rank Matrix Factorization}

\textbf{Classical SVD.} The Eckart--Young theorem (Eckart \& Young, 1936) states 
that the optimal rank-$r$ approximation of $W$ in Frobenius norm is the 
truncated SVD. Dobi-SVD uses this for initialization.

\textbf{Randomized SVD.} Halko et al. (2011) provide randomized algorithms. Not 
used here; exact SVD is feasible for our sizes.

\textbf{Symmetric vs. explicit factorization.} Standard implementations distribute 
$\Sigma^{1/2}$ between $U$ and $V$ (Golub \& Van Loan, 2013). Dobi-SVD 
uses symmetric form and pre-scales $V$ by the gate compensation factor.

\textbf{Weighted low-rank.} Golub et al. (1987); Srebro \& Jaakkola (2003). 
Activation-aware methods (Lin et al., 2023; Xiao et al., 2023) use 
activation statistics. Dobi-SVD uses unweighted SVD with \textbf{learnable gating}.

\subsection{Rank Selection Strategies}

\subsubsection{Fixed Rank}

Most methods use a single global $r$ (Yu et al., 2017). Simple but ignores 
layer heterogeneity.

\subsubsection{Energy-Based Rank}

Choose the smallest $r$ with cumulative spectral energy $\geq \eta$. 
Training-free. Used by MinimalSVD, MinimalNOVA, MinimalPRISM.

\subsubsection{Learned Rank (AdaLoRA)}

AdaLoRA (Zhang et al., 2023) learns per-layer rank budgets via 
importance scores. Related but different: AdaLoRA learns a \textbf{budget}, 
Dobi-SVD learns \textbf{which rank components survive}.

\subsubsection{Differentiable Pruning}

Louizos et al. (2018) use $L_0$ gates for structured pruning. Dobi-SVD 
applies a similar idea to \textbf{rank components} in SVD factorization.

\subsection{Parameter-Efficient Fine-Tuning}

\textbf{LoRA.} Hu et al. (2021) freeze base weights and train rank-$r$ adapters. 
Dobi-SVD trains LoRA + gates jointly.

\textbf{DoRA.} Liu et al. (2024) decompose weights into magnitude and direction. 
Dobi-SVD decomposes into rank components with gating.

\subsection{Attention Variants}

\textbf{Standard MHSA.} Vaswani et al. (2017).

\textbf{GQA.} Ainslie et al. (2023). Supported via \texttt{repeat\_kv}.

\subsection{Model Compression}

\begin{table}[H]
\centering
\small
\begin{tabular}{llll}
\toprule
\textbf{Method} & \textbf{Rank selection} & \textbf{Differentiable?} & \textbf{Post-hoc pruning?} \\
\midrule
DistilBERT & -- (distillation) & -- & -- \\
AWQ & -- (quantization) & -- & -- \\
SparseGPT & -- (pruning) & -- & -- \\
MinimalSVD & Energy &  &  \\
SVD-LLM V2 & Energy &  &  \\
MLA & Fixed latent &  &  \\
AdaLoRA & Learned budget &  & Partial \\
\textbf{Dobi-SVD} & \textbf{Learnable gates} & \textbf{ } & \textbf{ } \\
\bottomrule
\end{tabular}
\end{table}

\subsection{Position of Our Work}

Dobi-SVD differs from prior work along three axes:

\begin{table}[H]
\centering
\small
\begin{tabular}{lll}
\toprule
\textbf{Axis} & \textbf{Prior} & \textbf{Dobi-SVD} \\
\midrule
\textbf{Rank representation} & Integer & Continuous gates \\
\textbf{Differentiability} & No & Yes \\
\textbf{Post-hoc pruning} & Retrain & Threshold gates \\
\bottomrule
\end{tabular}
\end{table}

\textbf{Key distinction.} Unlike fixed or energy-selected rank, Dobi-SVD's rank 
is \textbf{learned jointly with recovery}. After training, gates provide a 
data-driven signal for which components to keep.
\begin{lstlisting}

---
\end{lstlisting}